
\pagebreak
\begin{center}
\textbf{Laboratory assignment 2: GraphQL.}
\end{center}


\medskip


\begin{enumerate}
\item Build a GraphQL API for your application (topic choosen in previous assignment).

\begin{itemize}
\item In Node.js you can use Apollo Server, in Asp.NET - HotChocolate nuget package.

\item Create GraphQL schema with defined types
\end{itemize}



\item Implement basic queries and mutations:

\begin{itemize}
\item Query all items

\item Query an item by ID

\item Add a new item mutation

\item Update item�s details via mutation

\item Delete an item
\end{itemize}




\item Provide query and mutation schema in a separate file.

\begin{itemize}
\item Example query:

\begin{lstlisting}
query {
  items(filter: "test", skip: 0, take: 5) {
    _id
    fieldOne
    fieldTwo
  }
}
\end{lstlisting}


\end{itemize}



\item Introduce query and mutation to your ui. You can use Apollo package if you use React, or just via plane POST requests of the format:

\begin{lstlisting}
const url = 'https://your-netlify-site.netlify.app/.netlify/functions/graphql';

const query = `
  query GetItems($filter: String) {
    items(filter: $filter) {
      _id
    fieldOne
    fieldTwo
    }
  }
`;

const variables = {
  filter: 'test', // Replace with the desired value or leave null for all values
};

fetch(url, {
  method: 'POST',
  headers: {
    'Content-Type': 'application/json',
  },
  body: JSON.stringify({
    query,
    variables,
  }),
})
  .then(response => response.json())
  .then(data => {
    console.log('Response data:', data);
  })
  .catch(error => {
    console.error('Error fetching data:', error);
  });
\end{lstlisting}


\end{enumerate}


